\section{Вступ та мотивація}

\subsection{Задача автоматизованої інспекції}

Сучасний логістичний склад є складним тривимірним простором: ряди металевих стелажів заввишки 2–6~м, вузькі коридори шириною 2–4~м, обмежені умови освітлення. Традиційні підходи до інспекції мають суттєві недоліки:

\begin{itemize}
  \item \textbf{Ручний обхід персоналом} — потребує значних людських ресурсів, є потенційно небезпечним в зонах з важкою технікою; не забезпечує повного покриття верхніх ярусів стелажів.
  \item \textbf{Стаціонарні камери} — фіксований кут зору, не дозволяють детально оглянути конструкції; висока вартість повного покриття великих приміщень.
\end{itemize}

БПЛА пропонують принципово інший підхід: рухома платформа здатна облітати весь об'єм приміщення, фіксувати відхилення від норми та в реальному часі передавати дані оператору. При об'єднанні кількох БПЛА у рій стає можливим паралельне покриття кількох зон, що скорочує час повної інспекції.

\subsection{Планування маршруту покриття}

\textbf{Boustrophedon Cellular Decomposition (BCD)} — класичний метод планування покриття, введений Choset~(2000)~\cite{Choset2000}. Ідея: розбити середовище на монотонні комірки, кожна з яких покривається рядковим обходом (з грецьк. βoustrophedon — «як віл оре»). У нашому сценарії кожен коридор між стелажами є монотонною коміркою, і дрон виконує один прямий прохід вздовж його осі.

\textbf{Методи на основі графів} (Spanning Tree Coverage, STC)~\cite{Hazon2005} гарантують оптимальне покриття для багатьох агентів, але потребують детальної попередньо відомої карти та централізованої координації. Для задачі з відомим планом складу та незалежними зонами BCD є достатнім і значно простішим.

\subsection{Алгоритми керування БПЛА}

\textbf{Геометричний SE(3)-трекер}: Lee, Leok \& McClamroch~(2010)~\cite{Lee2010} надали строгу геометричну основу для відстеження траєкторій на групі конфігурацій~SE(3), що гарантує майже-глобальну асимптотичну стійкість. Mellinger \& Kumar~(2011)~\cite{Mellinger2011} показали, що для квадрокоптера оптимальні траєкторії мінімізують похідну «snap» (четверта похідна позиції), а контролер формує сигнали керування через PD-регулятор позиції та SO(3)-регулятор орієнтації.

Pegasus Simulator~\cite{Pegasus2024} надає готову реалізацію \texttt{NonlinearController} на основі~\cite{Mellinger2011, Pinto2021}, яку ми використовуємо без модифікацій внутрішньої математики.

\subsection{Генерація траєкторій}

Lynch \& Park~(2017)~\cite{Lynch2017} у §9.2 доводять, що \textbf{квінтичний поліном}~(5-го ступеня) є мінімальним, що одночасно задовольняє граничним умовам по позиції, швидкості та прискоренню в обох кінцях. Нижчий ступінь (кубічний) залишає розрив прискорення, що призводить до неприпустимих стрибків моменту у квадрокоптері.

\subsection{IoT-протоколи для мобільних вузлів}

Стандарт MQTT~v5.0~\cite{MQTT5} розроблений для пристроїв з обмеженими ресурсами. Ключові переваги для нашого застосування:
\begin{itemize}
  \item фіксований службовий заголовок 2~байти (мінімальні витрати на повідомлення);
  \item механізм QoS~0/1/2 для різних рівнів критичності;
  \item топікова адресація без прямих з'єднань між вузлами — БПЛА не знає IP-адреси сервера моніторингу.
\end{itemize}

\subsection{Задача естафетної передачі при обмеженій батареї}

Квадрокоптер Iris в Pegasus має симульований ресурс батареї. При тривалих місіях необхідна автоматична ротація дронів. Ключова технічна вимога: дрон-наступник повинен знати \emph{точні світові координати} (X,~Y,~Z) місця, де попередник зупинився, щоб забезпечити відсутність пропусків у покритті.

Запропонований підхід відрізняється від класичного «handoff through planning»: замість попередньо розрахованого плану, позиція передається у реальному часі через той само IoT-канал (MQTT), який використовується для телеметрії.
